% migrate-to: execution-market-making, market-making-lob, market-structure-execution-quality, hft-market-design
\subsection{Price Discovery}

% migrated: Price Formation Process (PFP) -> def:price-formation-process; Informed Trading Price Discovery -> res:informed-trading-aggregation; Value-Trader Profit Source Segmentation -> def:outside-spread (Conditions, Pitfalls); Value Traders Supply Liquidity Outside the No-Trade Band -> def:outside-spread (Statement); No-Trade Is an Information Position -> def:outside-spread (Practice); Microstructure to Macro Prices, Volatility Scale Segmentation, Intraday and Overnight Returns, Correlation Interpretation -> def:intraday-overnight-returns (full text in notes/artefacts/market-making-lob/derivations.tex)

% assign: market-structure-execution-quality
\begin{mechanism} \hlt{Fragmentation and Price Formation}\\
Fragmentation can improve prices when venues remain synchronised by routing, arbitrage, and consolidated information. It can damage price formation when venues lose a common reference, when primary quotes become stale, or when market participants pause routing during stress.
\end{mechanism}

% assign: hft-market-design
\begin{mechanism} \hlt{Fast Liquidity and Price Discovery}\\
Fast market makers help price discovery when they update quotes and arbitrage stale venues. They hurt it when speed investment mainly earns queue priority or fee rebates without adding durable liquidity. Fee schedules therefore affect the price formation process, not only venue revenue.
\end{mechanism}

% assign: market-structure-execution-quality
\begin{caveat} \hlt{Outage Price Discovery}\\
If alternative venues keep trading while primary data or primary matching is unavailable, prices may temporarily split into parallel formation processes rather than one consolidated market price.
\end{caveat}
